\medskip
L'espace est rapporté à un repère orthonormé \Oijk.

On considère :

\begin{itemize}
\item les points C$(3~;~0~;~0)$, D$(0~;~2~;~0)$, H$(-6~;~2~;~2)$ et J$\left(\dfrac{-54}{13}~;~\dfrac{62}{13}~;~0\right)$;
\item le plan $P$ d'équation cartésienne $2 x+3 y+ 6 z-6 = 0$;
\item le plan $P'$ d'équation cartésienne $x - 2 y + 3 z - 3 = 0$;
\item la droite $(d)$ dont une représentation paramétrique est : $\begin{cases}x=- 8+\frac{1}{3} t \\ y = -1+\frac{1}{2} t \\ z=-4+t\end{cases}, \enskip t \in \R$
\end{itemize}

Pour chacune des affirmations suivantes, préciser si elle est vraie ou fausse, puis justifier la réponse donnée. Une réponse non argumentée ne sera pas prise en compte.

\medskip

\textbf{Affirmation 1 : } La droite $(d)$ est orthogonale au plan $P$ et coupe ce plan en H.

\medskip

\textbf{Affirmation 2 : } La mesure en degré de l'angle $\widehat{\mathrm{D C H}}$, arrondie à $10^{-1}$, est $17,3^{\circ}$.

\medskip

\textbf{Affirmation 3 : } Les plans $P$ et $P'$ sont sécants et leur intersection est la droite $\Delta$ dont une représentation paramétrique est : $\left\{\begin{array}{l c l}x&=&3-3 t \\ y&=&0 \\ z&=&\phantom{3- 3}t\end{array}\right.,\enskip t \in \R$.
\medskip

\textbf{Affirmation 4 : } Le point J est le projeté orthogonal du point H sur la droite (CD).

\medskip
